\begin{table}[htbp]
\centering
\scriptsize
\setlength{\tabcolsep}{4pt}
\begin{threeparttable}
\caption{Surveys, windows, sample sizes and occupation depth}
\label{tab:surveys}
\begin{tabular}{lllrrrrrrl}
\toprule
Country & Survey & Quarters & Employed & Min. & Cell & \% at 3+ & Occup. & \% cells & Tenure \\
 &  & (number) & per quarter & age & digits & digits & in cells & under floor &  \\
\midrule
Brazil$^{a}$ & PNADC & 2021Q1--2026Q2 (22) & 202\,100 & 14 & 3 & 99 & 124 & 42 & yes \\
Mexico$^{b}$ & ENOE & 2021Q1--2026Q2 (22) & 186\,700 & 15 & 2 & 61 & 43 & 20 & first int. \\
Colombia$^{c}$ & GEIH & 2022Q1--2026Q2 (18) & 91\,000 & 15 & 3 & 100 & 130 & 57 & yes \\
Argentina$^{d}$ & EPH & 2021Q1--2026Q1 (21) & 20\,700 & 10 & 2 & 0 & 40 & 62 & yes \\
Ecuador & ENEMDU & 2021Q3--2026Q1 (19) & 41\,700 & 15 & 3 & 100 & 130 & 72 & yes \\
Peru$^{e}$ & EPEN & 2022Q1--2026Q2 (18) & 48\,400 & 14 & 3 & 95 & 117 & 70 & no \\
Uruguay$^{f}$ & ECH & 2021Q3--2025Q4 (18) & 31\,100 & 14 & 3 & 100 & 126 & 73 & first int. \\
Bolivia$^{g}$ & ECE & 2021Q1--2025Q4 (20) & 27\,700 & 14 & 2 & 73 & 43 & 58 & no \\
South Africa$^{h}$ & QLFS & 2022Q1--2026Q2 (18) & 15\,600 & 15 & 3 & 94 & 118 & 85 & yes \\
Georgia$^{i}$ & LFS & 2021Q1--2025Q4 (20) & 6\,600 & 15 & 3 & 100 & 127 & 92 & no \\
Philippines$^{j}$ & LFS & 2021Q1--2025Q4 (19) & 122\,700 & 15 & 2 & 0 & 43 & 35 & no \\
\midrule
India$^{k}$ & PLFS & 2021Q1--2025Q4 (20) & 71\,500 & 15 & 3 & 99 & 127 & 76 & no \\
Nigeria$^{l}$ & NLFS & 2024Q1--2025Q2 (5) & 18\,900 & 15 & 3 & 98 & 130 & 86 & partial \\
\bottomrule
\end{tabular}
\begin{tablenotes}[flushleft]
\footnotesize
\item \textit{Notes:} Quarters gives the first and last quarter in the panel and, in parentheses, the number of quarters. PNADC is Brazil's Pesquisa Nacional por Amostra de Domic\'ilios Cont\'inua. Employed per quarter is the unweighted count of employed persons in the harmonised file, rounded to hundreds. Cell digits is the ISCO-08 level of the analysis cells (three where at least 90 percent of employment resolves to three digits, else two); \% at 3+ digits is the employment-weighted share of workers whose occupation code resolves to three or more digits. Occupations in cells and the share of cells under the 30-observation floor refer to the baseline year 2022. Tenure says whether the survey asks job tenure, which defines the new-hire share. India and Nigeria are outside the pooled estimates.
\item $^{a}$ Reweighted in August 2025 to the 2022 Census; one vintage kept.
\item $^{b}$ SINCO crosswalk reliable at two digits; tenure asked in first-interview quarters only.
\item $^{c}$ Starts at the 2022 redesign, so no 2021 pre-period.
\item $^{d}$ 31 urban agglomerations only; CNO 2017 crosswalk reliable at two digits.
\item $^{e}$ No tenure question; national file.
\item $^{f}$ Person-months (weight divided by three); tenure for first interviews only.
\item $^{g}$ COB 2009 follows ISCO in its first four digits; a quarter of codes stop at two.
\item $^{h}$ SASCO mapped through ISCO-88; the employment status question changed in 2025Q3.
\item $^{i}$ Small survey; 92 percent of three-digit cells fall under the floor.
\item $^{j}$ PSOC published at two digits; expanded rounds alternate with regular rounds; 2025Q3 missing.
\item $^{k}$ Annex only: 2025 redesign; first visits only before 2025Q2; NCO-2004 codes in 2021H1.
\item $^{l}$ Not in the estimates: no pre-period.
\end{tablenotes}
\end{threeparttable}
\end{table}
